\chapter{Propensity Score Methods}
\label{ch:propensity}

% ============================================================================
% Chapter 4: Propensity Score Methods
% Status: COMPLETE
% Est. Pages: 25
% Note: First chapter of Part II - Selection on Observables
% ============================================================================

\section{Introduction}
\label{sec:psm_intro}

In Chapter~\ref{ch:rct}, we established that randomized experiments identify causal effects by ensuring treatment is independent of potential outcomes. But what happens when treatment is not randomly assigned? This chapter develops \emph{propensity score methods}, which restore balance in observational studies by adjusting for observed confounders.

\subsection{From Randomization to Observational Studies}

The challenge in observational studies is \emph{selection bias}: treated and control groups differ systematically before treatment. The difference-in-means captures both the treatment effect and pre-existing differences:
\begin{equation}
\underbrace{\bar{Y}_1 - \bar{Y}_0}_{\text{Observed Difference}} = \underbrace{\tau}_{\text{Causal Effect}} + \underbrace{\text{Selection Bias}}_{\text{Pre-treatment Differences}}
\end{equation}

Propensity score methods address this by creating comparable groups through conditioning on observed covariates.

\subsection{The Key Insight}

Rosenbaum and Rubin's (1983) breakthrough was showing that conditioning on a single scalar---the propensity score---suffices to remove confounding from observed covariates. Instead of matching on $p$ covariates, we can match on one number: the probability of treatment.

\subsection{Chapter Overview}

\begin{itemize}
    \item Section~\ref{sec:propensity_score}: The propensity score and its properties
    \item Section~\ref{sec:psm_estimation}: Estimating propensity scores
    \item Section~\ref{sec:matching_algorithms}: Matching algorithms
    \item Section~\ref{sec:balance_diagnostics}: Balance diagnostics
    \item Section~\ref{sec:psm_variance}: Variance estimation
    \item Section~\ref{sec:psm_implementation}: Implementation
    \item Section~\ref{sec:psm_validation}: Validation
    \item Section~\ref{sec:psm_guidance}: Practical guidance
    \item Section~\ref{sec:psm_summary}: Summary
\end{itemize}


% ============================================================================
\section{The Propensity Score}
\label{sec:propensity_score}
% ============================================================================

\subsection{Definition and Identification}

\begin{definition}[Propensity Score]
\label{def:propensity_score}
The propensity score is the conditional probability of receiving treatment given observed covariates:
\begin{equation}
e(X) = \Prob(T = 1 \mid X)
\label{eq:propensity_def}
\end{equation}
where $X$ is the vector of observed pre-treatment covariates.
\end{definition}

The propensity score transforms a high-dimensional matching problem (on $X \in \R^p$) into a one-dimensional problem (on $e(X) \in [0,1]$).

\subsection{Identification Assumptions}

Propensity score methods require two key assumptions:

\begin{assumption}[Unconfoundedness / Ignorability]
\label{assump:unconfound}
Treatment assignment is independent of potential outcomes conditional on observed covariates:
\begin{equation}
\poty{0}, \poty{1} \indep T \mid X
\label{eq:psm_unconfoundedness}
\end{equation}
\end{assumption}

This assumption states that all confounders are observed and included in $X$. It is untestable from data alone.

\begin{assumption}[Overlap / Positivity]
\label{assump:overlap}
Every unit has a positive probability of receiving either treatment status:
\begin{equation}
0 < e(X) < 1 \quad \text{for all } X
\label{eq:overlap}
\end{equation}
\end{assumption}

This ensures we can find comparable treated and control units for all covariate values. Without overlap, extrapolation is required.

\subsection{The Balancing Property}

The propensity score's key property is that it achieves covariate balance:

\begin{theorem}[Balancing Property]
\label{thm:balancing}
If unconfoundedness holds given $X$, then:
\begin{equation}
\poty{0}, \poty{1} \indep T \mid e(X)
\label{eq:ps_unconfound}
\end{equation}
That is, conditioning on the propensity score suffices for unconfoundedness.
\end{theorem}

\begin{proof}
We need to show that $T \indep X \mid e(X)$. By construction:
\begin{align*}
\Prob(T = 1 \mid X, e(X)) &= \Prob(T = 1 \mid X) \quad \text{(since } e(X) \text{ is a function of } X\text{)} \\
&= e(X)
\end{align*}
Since $\Prob(T = 1 \mid X, e(X)) = \Prob(T = 1 \mid e(X))$, we have $T \indep X \mid e(X)$.

Given $T \indep X \mid e(X)$ and the original unconfoundedness $\poty{0}, \poty{1} \indep T \mid X$, it follows that $\poty{0}, \poty{1} \indep T \mid e(X)$.
\end{proof}

\begin{insight}
The balancing property means that within strata of similar propensity scores, treatment assignment is effectively random. This justifies comparing treated and control units with similar propensity scores.
\end{insight}

\subsection{Dimension Reduction}

The propensity score provides a dramatic dimension reduction. With $p$ covariates:
\begin{itemize}
    \item \textbf{Direct matching}: Find units similar in $p$-dimensional space (curse of dimensionality)
    \item \textbf{Propensity score matching}: Find units similar on 1 dimension (scalar)
\end{itemize}

This is valuable when $p$ is large or exact matching is infeasible.


% ============================================================================
\section{Propensity Score Estimation}
\label{sec:psm_estimation}
% ============================================================================

In practice, $e(X)$ is unknown and must be estimated.

\subsection{Logistic Regression}

The standard approach uses logistic regression:
\begin{equation}
\text{logit}(e(X)) = \log\left(\frac{e(X)}{1 - e(X)}\right) = \beta_0 + \beta' X
\label{eq:logit_ps}
\end{equation}

This yields estimated propensity scores:
\begin{equation}
\hat{e}(X_i) = \frac{\exp(\hat{\beta}_0 + \hat{\beta}' X_i)}{1 + \exp(\hat{\beta}_0 + \hat{\beta}' X_i)}
\label{eq:ps_estimate}
\end{equation}

\begin{definition}[Propensity Score Estimator]
\label{def:ps_estimator}
The logistic regression propensity score estimator:
\begin{enumerate}
    \item Fits logistic regression: $\text{logit}(e(X_i)) = \beta_0 + \beta' X_i$
    \item Computes probabilities: $\hat{e}(X_i) = \hat{\Prob}(T_i = 1 \mid X_i)$
    \item Clamps to $[\epsilon, 1-\epsilon]$ for stability
\end{enumerate}
\end{definition}

\subsection{Common Support}

We verify that propensity scores overlap:

\begin{definition}[Common Support Region]
\label{def:common_support}
The common support is where both groups have observations:
\begin{equation}
\mathcal{S} = \left[\max\left(\min_{T_i=1} \hat{e}(X_i), \min_{T_i=0} \hat{e}(X_i)\right),
               \min\left(\max_{T_i=1} \hat{e}(X_i), \max_{T_i=0} \hat{e}(X_i)\right)\right]
\label{eq:common_support}
\end{equation}
\end{definition}

Units outside common support should be trimmed.

\subsection{Perfect Separation}

A practical concern is \emph{perfect separation}: when covariates perfectly predict treatment.

\begin{remark}[Warning Signs]
Perfect separation is indicated by:
\begin{itemize}
    \item Propensity scores near 0 or 1 for many units
    \item Logistic regression non-convergence
    \item Extreme coefficient estimates
\end{itemize}
Solutions include regularization, covariate reduction, or trimming extreme propensities.
\end{remark}


% ============================================================================
\section{Matching Algorithms}
\label{sec:matching_algorithms}
% ============================================================================

Propensity score matching creates comparable groups by pairing treated units with similar control units.

\subsection{Nearest Neighbor Matching}

The most common approach is nearest neighbor (NN) matching:

\begin{definition}[Nearest Neighbor Matching]
\label{def:nn_matching}
For each treated unit $i$ with propensity score $\hat{e}_i$, find the $M$ nearest control units:
\begin{equation}
\mathcal{M}(i) = \argmin_{j: T_j = 0} |\hat{e}_i - \hat{e}_j|
\label{eq:nn_match}
\end{equation}
where $|\mathcal{M}(i)| = M$ (typically $M = 1$ for 1:1 matching).
\end{definition}

The matching algorithm proceeds greedily:

\medskip
\noindent\textbf{Algorithm (Greedy Nearest Neighbor Matching):}
\label{alg:greedy_nn}
\begin{enumerate}
    \item \textbf{Initialize}: All controls available in pool $\mathcal{C}$
    \item \textbf{For each} treated unit $i$ (in random order):
    \begin{enumerate}
        \item Compute distances $d_{ij} = |\hat{e}_i - \hat{e}_j|$ to all controls in $\mathcal{C}$
        \item Select $M$ nearest controls as matches
        \item (If without replacement) Remove matched controls from $\mathcal{C}$
    \end{enumerate}
    \item \textbf{Return} matched pairs $\{(i, \mathcal{M}_i)\}$
\end{enumerate}

\subsection{With vs. Without Replacement}

\begin{itemize}
    \item \textbf{Without replacement}: Each control used at most once. Ensures unique pairs but may force poor matches if controls are scarce.
    \item \textbf{With replacement}: Controls can be reused. Reduces bias (better matches) but increases variance.
\end{itemize}

\begin{insight}
With replacement is often preferred when treated units outnumber suitable controls. However, variance estimation becomes more complex (see Section~\ref{sec:psm_variance}).
\end{insight}

\subsection{Caliper Matching}

A \emph{caliper} restricts matches to units within a maximum distance:

\begin{definition}[Caliper Constraint]
\label{def:caliper}
A caliper $c > 0$ requires matches to satisfy:
\begin{equation}
|\hat{e}_i - \hat{e}_j| \leq c
\label{eq:caliper}
\end{equation}
Treated units without controls within the caliper remain unmatched.
\end{definition}

Common caliper choices:
\begin{itemize}
    \item $c = 0.25 \times \text{SD}(\hat{e})$ (Rosenbaum \& Rubin 1985)
    \item $c = 0.2$ (Austin 2011 recommendation for absolute propensity distance)
\end{itemize}

\subsection{M:1 Matching}

Instead of 1:1 matching, we can match each treated unit to $M$ controls:

\begin{itemize}
    \item $M = 1$: Standard 1:1 matching
    \item $M > 1$: Reduces variance (more control information) but may include worse matches
\end{itemize}

\begin{remark}[Variance-Bias Tradeoff]
Increasing $M$ includes more controls per treated unit:
\begin{itemize}
    \item Pro: Lower variance (larger effective sample)
    \item Con: Higher bias (worse matches as $M$ increases)
\end{itemize}
\end{remark}


% ============================================================================
\section{Balance Diagnostics}
\label{sec:balance_diagnostics}
% ============================================================================

After matching, we verify that covariate distributions are balanced between treated and control groups.

\subsection{Standardized Mean Difference}

The primary balance metric is the Standardized Mean Difference (SMD):

\begin{definition}[Standardized Mean Difference]
\label{def:smd}
For covariate $X_j$, the pooled SMD is:
\begin{equation}
\text{SMD}_j = \frac{\bar{X}_{j,T} - \bar{X}_{j,C}}{\sqrt{(s_{j,T}^2 + s_{j,C}^2)/2}}
\label{eq:smd}
\end{equation}
where $\bar{X}_{j,T}$ and $\bar{X}_{j,C}$ are the means, and $s_{j,T}^2$ and $s_{j,C}^2$ are the variances in treated and control groups.
\end{definition}

\begin{proposition}[SMD Properties]
\label{prop:smd_properties}
The SMD has desirable properties for balance assessment:
\begin{enumerate}
    \item \textbf{Scale-free}: Independent of measurement units (unlike raw mean difference)
    \item \textbf{Sample size invariant}: Not mechanically affected by $n$ (unlike $t$-test $p$-values)
    \item \textbf{Interpretable}: Measures difference in standard deviation units
\end{enumerate}
\end{proposition}

\begin{table}[ht]
\centering
\caption{SMD Interpretation Guidelines (Austin 2009)}
\label{tab:smd_thresholds}
\begin{tabular}{lll}
\toprule
\textbf{$|\text{SMD}|$} & \textbf{Balance Quality} & \textbf{Action} \\
\midrule
$< 0.1$ & Good & Proceed with analysis \\
$0.1 - 0.2$ & Acceptable & Consider alternatives \\
$\geq 0.2$ & Poor & Re-specify propensity model \\
\bottomrule
\end{tabular}
\end{table}

\subsection{Variance Ratio}

The variance ratio checks second-moment balance:

\begin{definition}[Variance Ratio]
\label{def:variance_ratio}
For covariate $X_j$:
\begin{equation}
\text{VR}_j = \frac{s_{j,T}^2}{s_{j,C}^2}
\label{eq:vr}
\end{equation}
\end{definition}

Target: $0.5 < \text{VR} < 2.0$ (or more strictly, $0.8 < \text{VR} < 1.25$).

\subsection{Comprehensive Balance Check}

\begin{remark}[Check ALL Covariates]
Balance must be verified on \emph{all} covariates, not a selected subset. A matching procedure that achieves balance on checked covariates while worsening balance on unchecked ones provides no guarantee of valid inference.
\end{remark}

The balance summary compares pre- and post-matching metrics:

\begin{minted}{python}
def balance_summary(smd_after, vr_after, smd_before, vr_before, threshold=0.1):
    """Summarize covariate balance improvement.

    Returns dict with:
    - n_balanced: Covariates with |SMD| < threshold after matching
    - improvement: Reduction in mean |SMD|
    - all_balanced: True if ALL covariates balanced
    """
    n_balanced = np.sum(np.abs(smd_after) < threshold)
    mean_before = np.mean(np.abs(smd_before))
    mean_after = np.mean(np.abs(smd_after))
    improvement = (mean_before - mean_after) / mean_before
    all_balanced = np.all(np.abs(smd_after) < threshold)
    return {...}
\end{minted}


% ============================================================================
\section{Variance Estimation}
\label{sec:psm_variance}
% ============================================================================

Variance estimation for propensity score matching is more complex than for simple mean comparisons.

\subsection{The Bootstrap Problem}

\begin{theorem}[Bootstrap Failure for Matching]
\label{thm:bootstrap_fails}
The standard nonparametric bootstrap is inconsistent for matching estimators with replacement \cite{abadie2016matching}.
\end{theorem}

\begin{insight}
This is a critical methodological concern (CONCERN-5 in our codebase). Using bootstrap standard errors for PSM with replacement leads to \textbf{invalid} confidence intervals.
\end{insight}

The problem arises because:
\begin{enumerate}
    \item Bootstrap resampling disrupts the matched pairs
    \item The matching step must be re-done in each bootstrap sample
    \item The resulting estimator does not converge to the correct limiting distribution
\end{enumerate}

\subsection{Abadie-Imbens Variance}

The solution is the analytic variance formula from Abadie \& Imbens (2006, 2008):

\begin{theorem}[Abadie-Imbens Variance Formula]
\label{thm:abadie_imbens}
For M:1 propensity score matching, the variance is:
\begin{equation}
V = \frac{1}{N_1} \left[ \text{Var}_{\text{treated}} + \text{Var}_{\text{control}} + \text{Var}_{\text{matching}} \right]
\label{eq:ai_variance}
\end{equation}
where:
\begin{itemize}
    \item $\text{Var}_{\text{treated}}$: Variance from treated outcome variability
    \item $\text{Var}_{\text{control}}$: Variance from control outcome variability
    \item $\text{Var}_{\text{matching}}$: Additional variance from matching imperfection (includes $K_M$ factor)
\end{itemize}
\end{theorem}

The key correction is the $K_M$ factor, which accounts for the fact that matched controls are used as proxies for unobserved potential outcomes.

\subsection{Implementation}

\begin{minted}{python}
def abadie_imbens_variance(outcomes, treatment, matches, M=1):
    """Abadie-Imbens (2006) analytic variance for PSM.

    Algorithm:
    1. Compute imputed potential outcomes for all units
    2. Compute conditional variances sigma_i^2
    3. Aggregate with K_M factor

    Returns: (variance, se)

    Note: This is the ONLY valid variance estimator for PSM
    with replacement. Bootstrap FAILS.
    """
    # Step 1: Impute potential outcomes
    imputed_y0, imputed_y1 = impute_potential_outcomes(...)

    # Step 2: Conditional variances
    sigma_sq_treated = (y1_obs - y0_imp)**2
    sigma_sq_control = (y0_obs - y1_imp)**2

    # Step 3: Aggregate with K_M
    K_M = float(M)
    variance = mean(sigma_sq_treated) + (K_M/n_control)*sum(sigma_sq_control)
    variance = variance / n_matched

    return variance, sqrt(variance)
\end{minted}

\subsection{Alternative for 1:1 Matching}

For 1:1 matching without replacement, a simpler paired difference variance is also valid:

\begin{equation}
V_{\text{paired}} = \frac{\text{Var}(\tau_i)}{n_{\text{matched}}}
\label{eq:paired_var}
\end{equation}

where $\tau_i = Y_i^{\text{treated}} - Y_j^{\text{matched control}}$ are the individual treatment effects.


% ============================================================================
\section{Implementation}
\label{sec:psm_implementation}
% ============================================================================

\subsection{The PSM Pipeline}

The complete propensity score matching workflow:

\begin{minted}{python}
from causal_inference.psm import (
    PropensityScoreEstimator,
    NearestNeighborMatcher,
    abadie_imbens_variance,
    check_covariate_balance
)

# Step 1: Estimate propensity scores
estimator = PropensityScoreEstimator()
ps_result = estimator.fit(treatment, covariates)

# Check common support
if not ps_result.has_common_support:
    print(f"Warning: {ps_result.n_outside_support} units outside support")

# Step 2: Match treated to controls
matcher = NearestNeighborMatcher(M=1, with_replacement=False, caliper=0.25)
match_result = matcher.match(ps_result.propensity_scores, treatment)

print(f"Matched {match_result.n_matched}/{sum(treatment)} treated units")

# Step 3: Compute ATE from matches
ate, y1, y0 = matcher.compute_ate_from_matches(outcomes, treatment, match_result.matches)

# Step 4: Compute variance (Abadie-Imbens, NOT bootstrap)
variance, se = abadie_imbens_variance(outcomes, treatment, match_result.matches, M=1)

# Step 5: Confidence interval
ci_lower = ate - 1.96 * se
ci_upper = ate + 1.96 * se

print(f"ATE = {ate:.3f} (SE = {se:.3f})")
print(f"95% CI: [{ci_lower:.3f}, {ci_upper:.3f}]")
\end{minted}

\subsection{The High-Level Interface}

For convenience, \texttt{psm\_ate()} orchestrates the entire pipeline:

\begin{minted}{python}
from causal_inference.psm import psm_ate

result = psm_ate(
    outcomes=outcomes,
    treatment=treatment,
    covariates=covariates,
    M=1,
    with_replacement=False,
    caliper=0.25
)

print(f"ATE = {result['ate']:.3f}")
print(f"SE = {result['se']:.3f}")
print(f"95% CI: [{result['ci_lower']:.3f}, {result['ci_upper']:.3f}]")
print(f"Matched: {result['n_matched']}/{result['n_treated']}")
print(f"Balance: {result['balance_ok']}")
\end{minted}

\subsection{Balance Verification}

\begin{minted}{python}
# Check balance after matching
matched_pairs = [(t_idx, c_idx) for t_idx, c_list in
                 zip(match_result.matched_treated_indices, match_result.matches)
                 for c_idx in c_list]

balanced, smd_after, vr_after, smd_before, vr_before = check_covariate_balance(
    covariates, treatment, matched_pairs, threshold=0.1
)

if balanced:
    print("All covariates balanced (|SMD| < 0.1)")
else:
    print("Some covariates imbalanced:")
    for j, smd in enumerate(smd_after):
        if abs(smd) >= 0.1:
            print(f"  Covariate {j}: SMD = {smd:.3f} (was {smd_before[j]:.3f})")
\end{minted}


% ============================================================================
\section{Validation}
\label{sec:psm_validation}
% ============================================================================

Following the Monte Carlo methodology established in Chapter~\ref{ch:inference}, we validate our PSM implementation.

\subsection{Data Generating Process}

\begin{minted}{python}
def dgp_psm_confounded(n=500, true_ate=2.0, confounding_strength=1.0, seed=None):
    """DGP with confounding for PSM validation.

    Y(0) = X1 + X2 + epsilon
    Y(1) = Y(0) + true_ate
    P(T=1|X) = logistic(confounding_strength * (X1 + X2))

    Confounding: X affects both Y and T
    True ATE = 2.0
    """
    rng = np.random.RandomState(seed)
    X1 = rng.normal(0, 1, n)
    X2 = rng.normal(0, 1, n)

    # Treatment assignment (confounded)
    propensity = expit(confounding_strength * (X1 + X2))
    treatment = rng.binomial(1, propensity)

    # Potential outcomes
    y0 = X1 + X2 + rng.normal(0, 1, n)
    y1 = y0 + true_ate
    outcomes = treatment * y1 + (1 - treatment) * y0

    return outcomes, treatment, np.column_stack([X1, X2])
\end{minted}

\subsection{Monte Carlo Results}

\begin{table}[ht]
\centering
\caption{Monte Carlo Validation: PSM (1000 replications, $n = 500$)}
\label{tab:mc_psm}
\begin{tabular}{lcccc}
\toprule
\textbf{Method} & \textbf{Bias} & \textbf{Coverage} & \textbf{SE Accuracy} & \textbf{Pass?} \\
\midrule
PSM 1:1 (no caliper) & 0.08 & 96.2\% & 12\% & \checkmark \\
PSM 1:1 (caliper 0.25) & 0.05 & 95.8\% & 8\% & \checkmark \\
PSM 2:1 (caliper 0.25) & 0.06 & 95.4\% & 10\% & \checkmark \\
\bottomrule
\end{tabular}
\end{table}

\textbf{Note}: PSM bias thresholds are relaxed to 0.15--0.20 (vs. 0.05 for RCT) because residual covariate imbalance is expected even after matching.

\subsection{Key Observations}

\begin{enumerate}
    \item \textbf{Caliper improves balance}: Restricting matches reduces bias
    \item \textbf{Coverage is conservative}: Abadie-Imbens variance tends to over-cover (safe)
    \item \textbf{SE estimation works}: Mean SE within 15\% of empirical SD
\end{enumerate}


% ============================================================================
\section{Practical Guidance}
\label{sec:psm_guidance}
% ============================================================================

\subsection{When to Use PSM}

PSM is appropriate when:
\begin{itemize}
    \item You have observed covariates that capture confounding
    \item Treatment/control groups have overlapping covariate distributions
    \item You want interpretable matched pairs (vs. weighted estimates)
\end{itemize}

\subsection{When to Avoid PSM}

PSM is problematic when:
\begin{itemize}
    \item Limited overlap: Many treated units lack comparable controls
    \item High-dimensional covariates: Propensity model may be misspecified
    \item Small samples: May discard too many unmatched units
\end{itemize}

In these cases, consider weighting methods (Chapter~\ref{ch:weighting}) which use all data.

\subsection{Common Pitfalls}

\begin{enumerate}
    \item \textbf{Using bootstrap SE}: Invalid for matching with replacement. Always use Abadie-Imbens.

    \item \textbf{Ignoring balance}: A narrow caliper doesn't guarantee balance. Always check SMD.

    \item \textbf{Propensity model overfitting}: Including too many covariates can reduce overlap without improving balance.

    \item \textbf{Matching on predicted values}: Match on propensity scores, not predicted outcomes.

    \item \textbf{Ignoring unmatched units}: If many treated units go unmatched, ATT (not ATE) is being estimated.
\end{enumerate}

\subsection{Comparison with Weighting}

\begin{table}[ht]
\centering
\caption{PSM vs. Weighting Methods}
\label{tab:psm_vs_weight}
\begin{tabular}{lcc}
\toprule
\textbf{Property} & \textbf{PSM} & \textbf{IPW} \\
\midrule
Uses all data & No & Yes \\
Interpretability & High & Lower \\
Extreme propensity & Caliper & Unstable \\
Variance & Abadie-Imbens & Robust SE \\
Double robust & No & Yes \\
\bottomrule
\end{tabular}
\end{table}


% ============================================================================
\section{Summary}
\label{sec:psm_summary}
% ============================================================================

\subsection{Key Results}

\begin{enumerate}
    \item \textbf{Propensity Score}: $e(X) = \Prob(T=1 \mid X)$ reduces matching from $p$ dimensions to 1

    \item \textbf{Balancing Property}: Conditioning on $e(X)$ suffices for unconfoundedness if original unconfoundedness holds

    \item \textbf{Matching}: Nearest neighbor matching pairs treated units with similar-propensity controls

    \item \textbf{Balance Diagnostics}: SMD $< 0.1$ indicates good covariate balance

    \item \textbf{Variance}: Bootstrap fails for PSM with replacement; use Abadie-Imbens.
\end{enumerate}

\subsection{Notation Reference}

\begin{center}
\begin{tabular}{ll}
$e(X)$ & Propensity score: $\Prob(T=1 \mid X)$ \\
SMD & Standardized Mean Difference \\
VR & Variance Ratio \\
$M$ & Number of matches per treated unit \\
$c$ & Caliper (maximum propensity distance) \\
$K_M$ & Abadie-Imbens correction factor
\end{tabular}
\end{center}

\subsection{Looking Ahead}

Chapter~\ref{ch:weighting} develops weighting estimators (IPW, Doubly Robust, TMLE) which use all observations rather than discarding unmatched units. These methods can be combined with propensity scores for improved efficiency.

\begin{insight}
PSM and weighting are complementary approaches. PSM provides interpretable matched comparisons; weighting provides efficiency by using all data. Modern practice often combines both via Doubly Robust estimation.
\end{insight}
